\subsection{Coarsest partitions for embedded smooth regularity}
\label{subsec:coarsest-embedded-smooth-partitions}

In this subsection $X\subset\mathbb R^N$ is an arbitrary semialgebraic
set, and $\Phi$ is a finite family of semialgebraic subsets of $X$.
Neither closedness nor local closedness of $X$ is assumed.  All closures
used to define partitions are relative to $X$.
For a semialgebraic set $A\subset\mathbb R^N$, let $\operatorname{Sm}(A)$
be the set of points at which $A$ agrees locally with an embedded
$C^\infty$ manifold without boundary, and put
\[
 \Sigma(A)=A\setminus\operatorname{Sm}(A).
\]
The local manifold dimension may vary on different components.
Thus $\operatorname{Sm}(A)$ uses embedded smoothness at every local
dimension.  In the semialgebraic setting these local manifolds are Nash
manifolds; see \cite[pp.~629--630]{Fernando2018}.

A finite semialgebraic partition is \emph{cellular} if every stratum is
locally closed in $X$ and its closure is a union of strata.  It is
\emph{$F$-regular} if $\Sigma(Y)$ is a union of strata for every union
$Y$ of strata that is closed in $X$.  This requirement includes all
closed unions, not just individual stratum closures.  Adaptation to
$\Phi$ means that every member of $\Phi$ is a union of strata.  We order
partitions by coarsening: $P$ is coarser than $Q$ when every stratum of
$Q$ is contained in a stratum of $P$.

The construction must control the singularities of every closed
union of strata at once. We first choose smooth pieces in decreasing
dimension, testing the finitely many relevant unions before choosing
the next pieces. To obtain coarseness, we intersect the Boolean
algebras of saturated sets. Stability under closure and the singular
locus makes this common coarsening admissible; minimizing the number
of strata then gives the coarsest partition.

We first record the exact classical semialgebraic input.

\begin{lemma}[The smooth locus in every local dimension]
\label{lem:all-dimensional-semialgebraic-smooth-locus}
For every semialgebraic $A\subset\mathbb R^N$, the set
$\operatorname{Sm}(A)$ is semialgebraic and relatively open in $A$.
Its complement $\Sigma(A)$ is relatively closed and nowhere dense.
If $A\ne\varnothing$, then $\dim\Sigma(A)<\dim A$.
These assertions also hold when $A$ is nonpure or nonclosed.
\end{lemma}

\begin{proof}
Stasica's theorem states separately for each integer $k$ that the set
$A^{(k)}$ of smooth points of dimension $k$ is semialgebraic
\cite[Theorem~4, p.~152]{Stasica2003Smooth}.  Consequently
\[
 \operatorname{Sm}(A)=\bigcup_{k=0}^N A^{(k)}
\]
is semialgebraic.  This is also the smooth-locus convention
$\operatorname{Sth}(A)$ and $\operatorname{Sth}_k(A)$ in
\cite[\S2.A, pp.~635--636]{Fernando2018}.  It is relatively open by
the local definition.  No global maximal-dimension restriction is
imposed on this union.

The same semialgebraicity conclusion follows from Koike--Shiota's
theorems on the nonsmooth points of fibres, applied to the constant
map $A\to\{0\}$
\cite[Theorems~2.1--2.2 and Corollary~2.4, p.~956]{KoikeShiota2007}.
Their Remark~0.1 on p.~953 identifies the semialgebraic
$C^\infty$ and analytic manifold conventions, and their argument
on pp.~961--962 includes sets that are not locally closed.

For the dimension assertion choose a finite pure-dimensional Nash
frontier partition of $A$.  Such partitions, adapted to any prescribed
finite family, are classical
\cite[Theorem~1.3 and the frontier refinement following it,
pp.~3--4]{Coste2007}.  Every stratum of maximal dimension is open in
$A$: the closure of any other stratum cannot meet it, by the frontier
condition and the strict drop of dimension on proper frontiers.
These maximal-dimensional strata lie in $\operatorname{Sm}(A)$.
Their complement has smaller dimension, so the same is true of
$\Sigma(A)$.

This global dimension drop also gives the required local density even
for a nonpure set.  If $\Sigma(A)$ had nonempty relative interior,
it would contain a nonempty relatively open semialgebraic subset $U$
of $A$, obtained by intersecting with a sufficiently small open ball.
Smoothness is local, so
$\Sigma(U)=U\cap\Sigma(A)=U$, contradicting the dimension drop
applied to $U$.  Thus $\Sigma(A)$ is nowhere dense.  The argument
applies to arbitrary semialgebraic subsets, so a set closed only in
$X$ is within its scope.
\end{proof}

\begin{theorem}[Coarsest smooth regular partition]
\label{thm:coarsest-embedded-smooth-partition}
There is a unique coarsest finite semialgebraic cellular $F$-regular
partition of $X$ adapted to $\Phi$.  Every one of its strata is
smooth.  The assertion makes no connectedness, local triviality, or
equisingularity requirement.
\end{theorem}

\begin{proof}
If $X$ is empty, its empty partition has all the asserted properties.
Suppose henceforth that $X$ is nonempty.  We first construct one
admissible partition and then prove the coarsest property.

Choose a finite pure-dimensional smooth frontier partition adapted
to $\Phi$.  We freeze pieces in decreasing order of dimension.  At
each stage keep a closed semialgebraic remainder $R\subset X$,
disjoint from the already frozen pieces, with a smooth frontier
partition adapted to $\Phi$, to the earlier pieces, and to every
closure in $X$ of a frozen piece.  Initially $R=X$.

Let $k=\dim R$, and let $S$ be one of the current $k$-dimensional
strata.  The frontier condition makes $S$ open in $R$.  All previously
frozen strata have dimension greater than $k$.  For every subset $J$
of the finite set of frozen strata, form the closed set
\[
 A_J=\bigcup_{T\in J}\overline T^{\,X}.
\]
The finitely many sets $\Sigma(A_J)\cap S$ are semialgebraic.
Remove all their relative boundaries in $S$.  Each of these
boundaries has dimension less than $k$, by the semialgebraic
dimension theorem in a $k$-manifold.  On the complement, each
membership pattern in the sets $\Sigma(A_J)\cap S$ is an open
semialgebraic subset of $S$.  Freeze every nonempty such piece.
These pieces are smooth of dimension $k$, and they are open in $R$.

Perform this operation simultaneously on all current $k$-strata.
The new remainder is closed in $R$, hence closed in $X$, and has
dimension less than $k$.  Choose on it a smooth frontier partition
adapted to the preceding partition and to all frozen closures.
The new remainder, being closed and disjoint from the frozen pieces,
cannot accumulate at a point of a frozen piece.  Within a given
$k$-stratum the newly frozen open pieces do not lie in each other's
closures.  Adaptation to their closures supplies all future frontier
relations.  The same statements hold between distinct $k$-strata by
the preceding frontier condition.  These observations preserve the
stated induction invariant.  The dimension decreases at each stage,
so after finitely many stages the frozen pieces give a finite smooth
frontier partition $P$ adapted to $\Phi$.

We verify the condition on every closed saturated union.  Let $S$ be
a final stratum of dimension $k$ and let $Y$ be a union of final
strata that is closed in $X$.  If $S$ is disjoint from $Y$, it is also
disjoint from $\Sigma(Y)$.  Suppose $S\subset Y$.  Collect the
strata $T\subset Y$ of dimension greater than $k$ for which
$S\subset\overline T^{\,X}$.  The frontier condition implies that
these are precisely the strata other than $S$ whose closures can
meet $S$.  At each point of $S$, a neighborhood avoids the finitely
many other stratum closures.  If the collection is empty, the germ
of $Y$ there is the smooth germ of $S$.  Otherwise the germ of $Y$
is the germ of the union $A_J$ of the collected closures: closedness
of $Y$ ensures that each of those closures lies in $Y$, and their
union already contains $S$.  All the collected strata had been
frozen before $S$ was constructed.  The membership test at that
stage therefore makes $\Sigma(A_J)\cap S$ either all of $S$ or
empty.  Locality of smoothness gives the same conclusion for
$\Sigma(Y)\cap S$.  Thus $P$ is $F$-regular.

For coarseness, associate to each admissible partition $P$ the finite
Boolean algebra $\mathcal B_P$ of all its saturated subsets.
Cellularity implies that this algebra is closed under closure in
$X$, because closure commutes with finite unions.  Its closed
members are also stable under $Y\mapsto\Sigma(Y)$, by
$F$-regularity.  If $P$ and $Q$ are adapted admissible partitions,
their intersection
\[
 \mathcal B=\mathcal B_P\cap\mathcal B_Q
\]
has the same properties and contains $\Phi$.  Its nonempty atoms
form a finite semialgebraic common coarsening $R$.

We check local closedness of those atoms.  For an atom $A$, both
$\overline A^{\,X}$ and
$F=\overline A^{\,X}\setminus A$ belong to $\mathcal B$, as does
$\overline F^{\,X}$.  If $A$ met $\overline F^{\,X}$, atomicity
would imply $A\subset\overline F^{\,X}$.  But the semialgebraic
frontier inequality gives
\[
 \dim\overline F^{\,X}=\dim F<\dim A,
\]
a contradiction; see \cite[Proposition~1.9, p.~5]{Coste2007}.
Hence $A=\overline A^{\,X}\setminus\overline F^{\,X}$ is open
in its closure.  Each atom closure is a union of atoms, so $R$ is
cellular.  The stability of $\mathcal B$ under $\Sigma$ on closed
members makes $R$ $F$-regular.

Moreover $\Sigma(\overline A^{\,X})\cap A$, being in
$\mathcal B$, is empty or all of $A$.  The latter is impossible:
$A$ is nonempty and open in its closure, while the nonsmooth locus
is nowhere dense there by
Lemma~\ref{lem:all-dimensional-semialgebraic-smooth-locus}.
Thus $A$ is open in the smooth locus of its closure, and is itself
an embedded smooth manifold, with local dimension allowed to vary.

Choose an adapted admissible partition $P$ with the smallest number
of strata; existence of such a minimum follows from the construction
and the well-ordering of the nonnegative integers.  The common
coarsening of $P$ and any other adapted admissible $Q$ just constructed
cannot have fewer strata than $P$.  Since it is coarser than $P$,
it must equal $P$.  Consequently $P$ is coarser than every $Q$.
Two partitions with this property coarsen each other and are equal,
which proves uniqueness.
\end{proof}

\begin{theorem}[Connected and pure-dimensional versions]
\label{thm:coarsest-smooth-partition-variants}
The preceding existence and uniqueness assertion holds separately
when each stratum is required to be connected, when each stratum
is required to have a single local dimension, and when both
requirements are imposed.  A connected smooth stratum automatically
has a single local dimension.
\end{theorem}

\begin{proof}
The construction in
Theorem~\ref{thm:coarsest-embedded-smooth-partition} already produces
pure-dimensional pieces.  For the connected version, freeze the
connected components of each open piece instead.  There are finitely
many such components and they are semialgebraic
\cite[Proposition~1.6, p.~5]{Coste2007}; they are open in the smooth
piece.  Refining the lower-dimensional remainder by their closures
preserves the same frontier and closed-remainder arguments.
Thus each of the stated classes is nonempty.

For connected strata, each connected component of a saturated set
is saturated: a connected stratum lies wholly in one component.
Consequently $\mathcal B_P$ is closed under taking connected
components, and so is $\mathcal B_P\cap\mathcal B_Q$ for two
connected partitions.  A disconnected atom would have a nonempty
proper connected component in that algebra, which contradicts
atomicity.  The common coarsening constructed above therefore has
connected strata.  Its smooth strata have locally constant manifold
dimension, so connectedness also makes that dimension constant.

For the class with pure-dimensional strata, use the local dimension sets
\[
 D_k(Y)=\{y\in Y:\dim_yY\ge k\}
\]
for each closed saturated $Y$.  For a pure-dimensional frontier
partition $P$, one has the exact formula
\[
 D_k(Y)=
 \bigcup_{\substack{T\in P,\ T\subset Y\\\dim T\ge k}}
       \overline T^{\,X}.
\]
Indeed the local dimension at a point of a stratum $S\subset Y$
is the maximum of the dimensions of the strata $T\subset Y$ whose
closures contain $S$.  A neighborhood avoids every other stratum,
and every neighborhood meeting such a pure-dimensional $T$ meets
an open subset of $T$ of its full dimension.  The displayed finite
union is closed in $X$ and saturated.  It shows directly that each
$\mathcal B_P$, and hence the intersection algebra for two pure
partitions, is stable under $D_k$ on closed members.

For an atom $A$ of that intersection, each set
$A\cap D_k(\overline A^{\,X})$ is empty or all of $A$.
The earlier proof makes $A$ open in the smooth locus of its closure,
so the local dimension of $\overline A^{\,X}$ on $A$ is the
manifold dimension of $A$.  The preceding all-or-none tests for all
$k$ force that dimension to be constant.  Thus the common coarsening
also belongs to the pure-dimensional class.  Applying the same
minimum-number-of-atoms argument in each class proves the assertions.
\end{proof}

